\documentclass[11pt]{article}

\newcommand{\cnum}{Math 245B}
\newcommand{\ced}{Winter 2026}
\newcommand{\ctitle}[4]{\title{\vspace{-0.5in}\cnum, \ced\\Assignment #1: #2}\author{\vspace{-0.35in}\\#3, #4}}
\usepackage{enumitem}
\usepackage[usenames,dvipsnames,svgnames,table,hyperref]{xcolor}
\usepackage{amsmath, amsfonts}
\usepackage{graphicx} % Required for inserting images
\usepackage{float}
\usepackage{listings}
\usepackage{xcolor}
\usepackage{hyperref}
\usepackage{comment}

\renewcommand*{\theenumi}{\alph{enumi}}
\renewcommand*\labelenumi{(\theenumi)}
\renewcommand*{\theenumii}{\roman{enumii}}
\renewcommand*\labelenumii{\theenumii.}
\newcommand{\notiff}{%
  \mathrel{{\ooalign{\hidewidth$\not\phantom{"}$\hidewidth\cr$\iff$}}}}

\definecolor{codegreen}{rgb}{0,0.6,0}
\definecolor{codegray}{rgb}{0.5,0.5,0.5}
\definecolor{codepurple}{rgb}{0.58,0,0.82}
\definecolor{backcolour}{rgb}{0.95,0.95,0.92}
\definecolor{header}{HTML}{205459}
\definecolor{solution}{HTML}{204d52}

\newcommand{\solution}[1]{{{\textcolor{header}{Solution:} \textcolor{solution}{#1}}}}
\setlist[enumerate]{label=(\alph*)}

\lstdefinestyle{mystyle}{
    backgroundcolor=\color{backcolour},
    commentstyle=\color{codegreen},
    keywordstyle=\color{magenta},
    numberstyle=\tiny\color{codegray},
    stringstyle=\color{codepurple},
    basicstyle=\ttfamily\footnotesize,
    breakatwhitespace=false,
    breaklines=true,
    captionpos=b,
    keepspaces=true,
    numbers=left,
    numbersep=5pt,
    showspaces=false,
    showstringspaces=false,
    showtabs=false,
    tabsize=2
}


\begin{document}
\ctitle{\#1}{}{Asher Christian}{006-150-286}
\date{}
\maketitle

\section{Monday}
Lemma: $\{X_n\}$ is a martingale w.r.t  $ \{ \mathcal{F}_n\}$ there exists $X_\infty$ such that
\[
    X_n = \mathbb{E}[X_\infty | \mathcal{F}_n] \iff \{X_n\} \text{ is u.i}
\] 
Proof: $\implies$  $\{ \mathbb{E}[ X | G] : G \in \mathcal{F}$\} is Uniformly integrable.\\
Proof of $\impliedby$  $\{X_n\}$ is uniformly integrable  $\implies$ 
 \[
\sup_n \mathbb{E}|X_n| < \infty
\] 
By martingale convergence theorem there exists $X_\infty$ such that $X_n \to X_\infty$ almost sure and $L_1$ since it is uniformly integrable.
Need to prove that for all $n$
\[
    \mathbb{E}[X_\infty | \mathcal{F}_n] = X_n
\] 
$\iff$ for all  $A \in \mathcal{F}_n$
\[
    \mathbb{E}[X_n 1_A ] = \mathbb{E}[X_\infty 1_A]
\] 
We know
\[
    \mathbb{E}[X_{n+1} | \mathcal{F}_n] = X_n \qquad \mathbb{E} X_{n+1} 1_A = \mathbb{E}[X_n 1_A]
\] 
for any $m > n$
 \[
     \mathbb{E}[X_m 1_A] = \mathbb{E}[X_n 1_A]
\] 
so
\[
    \mathbb{E}[X_n 1_A] = \lim_{m\to \infty}  \mathbb{E}[X_m 1_A] = \mathbb{E} \lim_{m\to \infty} X_n 1_A = \mathbb{E}[X_\infty 1_A]
\] 
since the convergence is $L_1$\\
Theorem: Levy 0-1 law.
\[
\{\mathcal{F}_n\} \quad F_\infty = \sigma(\bigcup \mathcal{F}_n) \quad A \in \mathcal{F}_\infty
\] 
filtration. Then
\[
    \mathbb{P}(A | \mathcal{F}_n) =  \mathbb{E}[1_A | \mathcal{F}_n] \to 1_A \text{a.s.}
\] 
Proof choose $X = 1_A$ and use levy's martingale convergence theorem.
Applicagtion: Kolomogorov 0-1 law.
If  $\{X_n\}$ are independent random variables. and $ \mathcal{F}_n = \sigma(X_1,\dots,X_n)$ and $\tau_N = \sigma( \mathcal{F}_{n+1} \dots )$ then $\tau = \bigcap \tau_n$ and  $A \in \tau \implies \mathbb{P}(A) \in \{0,1\}$ 
Use Levy's 0-1 law we have
\[
    \mathbb{E}[1_A | \mathcal{F}_n] \to 1_A
\] 
but $1_a \in \tau \subset \tau_n$ so $1_A \perp \mathcal{F}_n$ so
\[
    \mathbb{E}[1_a | \mathcal{F}_n] = \mathbb{E}[1_A] = \mathbb{P}(A)
\] 
and if
\[
\mathbb{P}(A) \to 1_A
\] 
then $1_A$ must be one or zero everyhwere and so
 \[
     \mathbb{E}[1_A] = \mathbb{P}(A) \in \{0,1\}
\] 
Thoerem (Conditional Dominated Convergence Theorem):\\
$ \{ \mathcal{F}_n\} \nearrow$ and $X_n \to X_\infty$ almost sure and 
\[
|X_n| < Z \quad a.s.
\] 
for all $n$. with $ \mathbb{E}|z| < \infty$
then
dct says
 \[
     \mathbb{E}[X_n] \to \mathbb{E}[X_\infty]
\] 
and we say
\[
    \mathbb{E}[X_n | \mathcal{F}_n] \to \mathbb{E}[X_\infty | \mathcal{F}_\infty]  \quad \text{a.s.}
\] 
Note we already learned that
\[
    \mathbb{E}[X_n | \mathcal{F}] \to \mathbb{E}[X_\infty | \mathcal{F}]
\] 
but that is different.\\
Proof:
\[
    \mathbb{E}[X_n | \mathcal{F}_n]  - \mathbb{E}[X_\infty | \mathcal{F}_\infty] 
    = \mathbb{E}[X_n | \mathcal{F}_n] + \mathbb{E}[X_\infty | \mathcal{F}_n] - ( \mathbb{E}[X_\infty | \mathcal{F}_\infty] - \mathbb{E}X_\infty | \mathcal{F}_n])
\] 
by Levy upward convergennce theorem the second term goes to zero so it suffices to check
\[
 |   \mathbb{E}[X_n | \mathcal{F}_n] - \mathbb{E}[X_\infty | \mathcal{F}_n] |\le ?
\] 
so we define
\[
    V_k = \sup_{n > k} |X_n - X_\infty|
\] 

\section{Wednesday}
Need to prove for $X_n \to X_\infty $ almost sure, and $ \mathcal{F}_n \nearrow \mathcal{F}_\infty$ and $|X_n| \le Z$ for all  $n$ that
 \[
     \lim_{n\to \infty} |\mathbb{E}[X_n | \mathcal{F}_n] -  \mathbb{E}[X_\infty | \mathcal{F}_\infty]| = 0
\] 
We want 
\[
    | \mathbb{E}(X_n | \mathcal{F}_n) - \mathbb{E}[X_\infty | \mathcal{F}_n] | \to 0
\] 
Define $V_k = \sup_{n \ge k} | X_n - X_\infty|$.
\[
    lhs \le \mathbb{E}[V_k | \mathcal{F}_n] \qquad n \ge k
\] 
we have
\[
    \lim_{n\to \infty} \mathbb{E}[V_k | \mathcal{F}_n] = \mathbb{E}[ V_k | \mathcal{F}_\infty]
\] 
by levy martingale convergence theorem.
We need to check that $ \mathbb{E}[|V_k|] < \infty$ is true since $ |V_k| \le 2Z$ is integrable.
we have
\[
    \lim_{n\to \infty} | \mathbb{E}(X_n | \mathcal{F}_n) - \mathbb{E}(X_\infty | \mathcal{F}_n)| \le \mathbb{E}[V_k | \mathcal{F}_\infty]
\] 
for all $k$ so it is correct for the infimum.  $X_n \to X_\infty$ almost sure implies $V_k \to 0$ almost sure and  $|V_k| \le 2Z$ and  $E Z < \infty$ so with dominated convergence theorem
\[
    \inf_k \mathbb{E}[V_k | \mathcal{F}_\infty] = \lim_{k\to \infty} \mathbb{E}[V_k | \mathcal{F}_\infty]  = 0
\] 
so
\[
\lim_{n\to \infty} | \mathbb{E}(X_n | \mathcal{F}_n) - \mathbb{E}(X_\infty | \mathcal{F}_n)| = 0
\] 

Levy's backwards martingale.
Normally we have $ \mathcal{F}_1, \mathcal{F}2,\dots \mathcal{F}_n, \dots \mathcal{F}_\infty$.
Backwards we have
\[
   \dots \mathcal{F}_{-3}, \mathcal{F}_{-2}, \mathcal{F}_{-1}, \mathcal{F}_0 
\] 
\[
    \{ \mathcal{F}_n\}_{n < 0} \qquad \mathcal{F}_n \subset \mathcal{F}_m \qquad n \le m
\] 
\[
    \mathcal{F}_{-\infty} = \bigcap_{n < 0} \mathcal{F}_n
\] 
is a sigma algebra.
Backward martingale $X_n \in \mathcal{F}_n$ and \\
$ \mathbb{E}[X_{n+1} | \mathcal{F}_n] = X_n$ \\
Result:
\emph{
    Levy BW martingale Convergence theorem\\
    \[
        X_{-n} \to X_{-\infty} \qquad \text{ a.s. and $L_1$}
    \] 
}
renomralization group method.
We skip the proof since $X_{-\infty} = \mathbb{E}[X_{-1} | \mathcal{F}_\infty]$ since we use 
the same argument as forward process and the $X_n$ are uniformly integrable so a.s. implies  $L_1$ \\
Applications:
Prove strong law of large numbers. We have $X_1,\dots,X_n$ i.id and $E |X_i| < \infty$
Define 
\[
    \mathcal{F}_{-n} := \mathcal{E}_n \qquad \mathcal{E}_n \subset \mathcal{F}, A \in \mathcal{E}_n, \pi \in S_n \pi(A) = A
\] 
symmetric sigma algebra with respect to the first $n$.
H-S 01 law.
We have
\[
\dots \subset \mathcal{E}_3 \subset \mathcal{E}_2
\] 
Let for $n > 0$
\[
    Y_{-n} = \mathbb{E}[X_1 | \mathcal{E}_n]
\] 
\[
    Y_{-2} + \mathbb{E}[X_1 | \mathcal{E}_2] = \frac{X_1 + X_2}{2}
\] 
since
\[
    \mathbb{E}[X_1 | \mathcal{E}_2] = \mathbb{E}[X_2 | \mathcal{E}_2] = \frac{1}{2} \mathbb{E}[X_1 + X_2 | \mathcal{E}_2]  = \frac{X_1 + X_2}{2}
\] 
since the last one is $ \mathcal{E}_2$ measurable.
So
\[
    Y_{-n} = \frac{\sum_{k=1}^{n}X_k}{n} \to Y_{-\infty} 
\] 
exists

\section{Friday}
SLLN
\[
\lim_{n\to \infty} \frac{S_n}{n} \to Z \qquad a.s.
\] 
and $Z$ is a random variable.
\[
    \mathbb{E}[Z] = \mathbb{E}[X_1]
\] 
since
\[
    Z = \lim_{n\to \infty } \mathbb{E}[X_1 |  \mathcal{E}_n] 
\] 
almost sure and $L_1$ so
\[
    \mathbb{E}[Z] = \mathbb{E}[X_1] 
\] 
Need to prove that $Z$ is constant\\
First 
\[
    \lim_{n\to \infty} \frac{X_1 + \dots + X_n}{n } = \lim_{n\to \infty} \frac{X_2 + \dots + X_{n+1}}{n} = Z^{(1)}
\] 
and 
\[
    Z = Z^{(1)} \quad \text{a.s.}
\]
and
\[
Z^{(1)} \perp X_1 \implies Z \perp X_1
\] 
We have
\[
    Z^{(k)} = Z \text{ a.s.}
\] 
and
\[
Z^{(k)} \perp \sigma(X_1,\dots,X_k) = \mathcal{F}_k
\] 
so
\[
Z \perp \mathcal{F}_k
\]
and
\[
z \perp \bigcup_{k} \mathcal{F}_k \implies Z \perp \sigma( \bigcup_{k} \mathcal{F}_k) = \mathcal{F}_\infty
\] 
but $Z \in \mathcal{F}_\infty$. If $X \in \mathcal{F} \qquad X \perp \mathcal{F}$ then $X$ is a constant.
We have
\[
    \mathbb{P}(X \le a) \in \{0,1\}
\] 
for all $a$ thus  $X$ is constant.\\


\end{document}
